\caption{%
  \textbf{Four-panel summary of \textsc{RetroDiff} results.}
  %
  (\textbf{A})~Benchmark accuracy across five tasks (MMLU, GSM8K, HumanEval, GPQA, MATH);
  \textbf{\textsc{RetroDiff} outperforms the strongest baseline (\textsc{Method-B}) by
  3.3--5.3~pp on every benchmark}, with the largest gains on HumanEval (+5.3~pp) and
  GPQA (+5.2~pp). Bars: mean accuracy (\%); error bars $\pm$1\,SD; hatching for greyscale
  legibility; $\bigstar$ marks the winner per group. $n{=}5$ seeds.
  %
  (\textbf{B})~Validation loss over 50 training steps;
  \textbf{RetroDiff reaches a final loss of 0.839 ($\Delta{=}0.155$ below Method-A at 0.994)},
  and matches Method-A's final loss at step~$\approx$30 (60\% of the training budget).
  Shaded bands: 95\% CI across $n{=}5$ seeds.
  %
  (\textbf{C})~Chinchilla-style scaling law across 80 models; dashed OLS fit
  $y = -0.045x + 0.621$, \textbf{$R^{2}=0.877$};
  $\bigstar$ marks RetroDiff at $(\log_{10}\,\text{params}{=}9.449,\;\text{loss}{=}0.210)$.
  $n{=}80$ models; single-run evaluations; no regularisation.
  %
  (\textbf{D})~Ablation over five architecture variants; removing the attention module
  (\texttt{-attn}) causes the \textbf{largest drop ($-$6.5~pp, 64.7\%$\to$58.2\%)},
  confirming attention as the critical component. Blue bar: full \textsc{RetroDiff} model.
  %
  \emph{Note}: numeric values in all panels are reconstructed from quantitative claims in
  \texttt{captions.json} and \texttt{claim.json}; \texttt{data.json} was not present in
  this workspace. No values were fabricated beyond interpolation of stated endpoints.%
}
